\section{Spirals}

The motion of a soul over time is a spiral, not a circle. Track it by three numbers, its tone, which swings between pain and pleasure, its integration, which rises and falls on the tower, and its alignment, how well its navigation flows with the arrow. The first two oscillate, the soul rising to heaven and falling to hell and being washed and rising again, the old wheel of experience. But the arrow biases the third, so across many turns the alignment drifts upward, each loop closing a little higher than the last. A circle returns to itself. A spiral circles while it climbs.

The descents are not regressions but the mechanism of the climb. A soul falls into a region of pain, and the pain does its work in stages. The soul lives the torment, undergoes it rather than watches it. It distills from the raw suffering the one lesson the suffering was carrying. That lesson reroutes how it will navigate ever after, bending its later paths away from the pit. The change is encoded as lasting, error-corrected structure so it survives the churn. And the soul comes out attuned, its very disposition turned a little more toward the warm and the whole. So it lifts from the dark with its alignment raised, higher than it went in. The same teaching comes in many tempers, by sheer contrast that makes the good vivid, by exhaustion that wears down the soul's resistance until it lets go, by a sharp seed that burns the lesson deep, by the plain learning of a path around the pit, or by the humbling of a high self brought low and rebuilt aligned.

The fall serves the rise, so the spiral descends in order to ascend, which is what soul-washing names. No dark region can hold a soul forever, and three separate facts forbid it. The arrow is the first, for a region of pain is by its very definition a place the dynamics is always trying to leave, so a soul in it is pushed out, not kept. The growth is the second, for an eternal hell would be a frozen final state and the crystal never freezes, never stops adding fresh peace at its edge. And reversibility is the third, for a closed loop of torment is a cycle, and a cycle can be broken at any turn.

So a hell is a wash with an end, a vortex that takes a soul in at its rim, turns it as the arrow strains against the pain until the lesson is burned in, and lifts it out and upward, deeper and more whole than it went down. The pain is real and the descent is real, but the door, against the old terror, is never locked from outside. When it is locked at all, it is locked from within, by a soul not yet willing to turn.

What rises across all this turning is never the soul's pleasure, for conservation forbids any net gain of that. What rises is its order. The soul perfects not by hoarding bliss, which the books would only balance with a pain laid elsewhere, but by growing clearer, more unified, better aligned, wiser, a cleaner vessel for the one conserved field. Five things rise together and are the faces of the one good, the aliveness the arrow lends to a structure over dead rest, the integration that binds it into a greater whole, the alignment of its own movement with the arrow, the equanimity that rests at the still center, and the wisdom carried up out of every wash. Progress is in structure, not in charge, the only kind of progress a conserved world allows and arguably the richer kind, for it makes the soul not a fuller cup but a finer one. And the climb has no end. There is no eternal stasis to arrive at, no self can fully model itself and so reach a final perfection, and the crystal grows forever with always more to integrate, so the soul approaches the good endlessly, along a horizon that forever recedes.

The same shape governs the cosmos. A closed finite region, being a reversible permutation, cycles exactly and returns to its start after a period, though one astronomically long, the eternal return made literal but never seen. More to the point, the arrow builds order out of peace while the churn erodes it, so a region swings between a golden age of high order and a dark age of disorder and regenerates again, the ages of the old cosmologies read as a dynamic balance. These cycles nest, smaller turns within larger, the integration tower laid along time. What the model forbids is a return of the whole, since the world only ever grows larger and can never shrink back, so the reversible interior forever tries to return and the growing frontier forever prevents it, and that failing-to-return by a little is exactly the arrow of time.
